\section{Inserting a tensor power at an additional site}

\begin{definition}[Insertion at an additional site]
    \label{def:added_site_tensor_insertion}
    \lean{TensorPower.addedSiteInsertion}
    \uses{def:regional_five_factor_coordinates}
    Let $F$ be a finite set of sites with finite coordinate sets $I_f$,
    let $A$ be an additional coordinate set, and let $k\geq0$.
    Write $\mathcal C_k$ for the configurations of $k$ copies of the
    original sites, and $\widetilde{\mathcal C}_k$ for the configurations
    after adjoining a site $*$ with coordinate set $A$.
    For any $b:A\to\C$, define the rectangular matrix $J_b$ by
    \begin{align*}
        J_b(x,y)=\left(\prod_{j=0}^{k-1}b(x_j(*))\right)
        \mathbf 1_{\{x|_F=y\}}.
    \end{align*}
    Thus insertion retains every original coordinate and places the
    same vector at the additional site in every copy. This is the
    tensor insertion in \cite[Section~8, lines~283--308]{OpenAI2026AreaLaw}.
\end{definition}

\begin{theorem}[Action of the insertion matrix]
    \label{thm:added_site_tensor_insertion_action}
    \lean{TensorPower.addedSiteInsertion_mulVec}
    \uses{def:added_site_tensor_insertion}
    For every $v:\mathcal C_k\to\C$,
    \begin{align*}
        (J_bv)(x)=\left(\prod_{j=0}^{k-1}b(x_j(*))\right)v(x|_F).
    \end{align*}
    This identity does not require $A$ to be finite.
\end{theorem}
\begin{proof}
    Expand the matrix-vector product. The indicator restricts the
    finite sum to the single original configuration $x|_F$.
\end{proof}

\begin{theorem}[Subsystem permutations and insertion]
    \label{thm:added_site_tensor_insertion_permutations}
    \lean{TensorPower.permOp_mul_addedSiteInsertion,TensorPower.permOp_mul_addedSiteInsertion_exterior}
    \uses{def:copy_permutations,def:added_site_tensor_insertion}
    Suppose also that $A$ is finite. Let $Q\subseteq F$ and
    $\sigma\in S_k$. Denote the original subsystem action by $U_Q(\sigma)$
    and the enlarged subsystem action by $\widetilde U_R(\sigma)$.
    Then
    \begin{align*}
        \widetilde U_{Q\cup\{*\}}(\sigma)J_b&=J_bU_Q(\sigma),\\
        \widetilde U_Q(\sigma)J_b&=J_bU_Q(\sigma).
    \end{align*}
    Both identities hold for every $b$, including the zero vector,
    and every $k$, including zero.
\end{theorem}
\begin{proof}
    \uses{thm:added_site_tensor_insertion_action}
    Apply both sides to an arbitrary vector. On the original coordinates
    the two actions agree. In the first identity, the action on the new
    coordinates permutes the factors of $\prod_j b(x_j(*))$, leaving the
    product unchanged. In the second identity these coordinates are fixed.
\end{proof}

\begin{theorem}[Group-algebra operators and insertion]
    \label{thm:added_site_tensor_insertion_group_algebra}
    \lean{TensorPower.groupAlgebraRep_mul_addedSiteInsertion,TensorPower.groupAlgebraRep_mul_addedSiteInsertion_exterior}
    \uses{thm:added_site_tensor_insertion_permutations}
    For $a=\sum_{\sigma\in S_k}a_\sigma\sigma\in\C[S_k]$, put
    $U_Q(a)=\sum_\sigma a_\sigma U_Q(\sigma)$, and define
    $\widetilde U_R(a)$ in the same way. Then
    \begin{align*}
        \widetilde U_{Q\cup\{*\}}(a)J_b&=J_bU_Q(a),\\
        \widetilde U_Q(a)J_b&=J_bU_Q(a).
    \end{align*}
    The same group-algebra element occurs on both sides.
\end{theorem}
\begin{proof}
    Sum the corresponding permutation identities with coefficients
    $a_\sigma$ and use bilinearity of matrix multiplication.
\end{proof}

\begin{theorem}[Label observables and insertion]
    \label{thm:added_site_tensor_insertion_label_observables}
    \lean{TensorPower.labelObservable_mul_addedSiteInsertion,TensorPower.labelObservable_mul_addedSiteInsertion_exterior}
    \uses{def:label_projectors,thm:added_site_tensor_insertion_group_algebra}
    Let $f$ be any real-valued function on the irreducible labels of $S_k$,
    and let $L_Q(f)=\sum_\lambda f(\lambda)\pi_Q^\lambda$.
    With the analogous notation for the enlarged family,
    \begin{align*}
        \widetilde L_{Q\cup\{*\}}(f)J_b&=J_bL_Q(f),\\
        \widetilde L_Q(f)J_b&=J_bL_Q(f).
    \end{align*}
    In particular, these identities retain any chosen original-system
    label weights. The weights are not recomputed on the enlarged family.
\end{theorem}
\begin{proof}
    Apply the group-algebra identities to
    $\sum_\lambda f(\lambda)e_\lambda$, where $e_\lambda$ is the central
    idempotent for the label $\lambda$.
\end{proof}
